\chapter{La suma: primeros pasos}\label{ch:g1:addition}

Tres canicas rojas y dos canicas azules: juntas, ¿cuántas son? Juntar
y contar el conjunto es \emph{sumar} --- la primera
operación, con sus dos signos famosos $+$ y $=$.

\section{Juntar}

\begin{definition}[La suma]\label{def:g1:addition:def}
\emph{Sumar} es juntar y contarlo todo. Escribimos
\[
3 + 2 = 5
\]
y se lee: «tres más dos es igual a cinco». El resultado, $5$, es la
\emph{suma}\index{suma} de $3$ y $2$.
\end{definition}

\begin{omfigure}
\begin{tikzpicture}[scale=0.6]
  \foreach \i in {0,1,2} \fill[omThm] (\i*1.0,0) circle (0.3);
  \node[font=\large] at (3.0,0) {$+$};
  \foreach \i in {0,1} \fill[omDef] (3.8+\i*1.0,0) circle (0.3);
  \node[font=\large] at (5.8,0) {$=$};
  \foreach \i in {0,1,2} \fill[omThm] (6.6+\i*1.0,0) circle (0.3);
  \foreach \i in {0,1} \fill[omDef] (9.6+\i*1.0,0) circle (0.3);
  \node[below, font=\small] at (5.6,-0.7)
    {$3 + 2 = 5$: tres y dos son cinco};
\end{tikzpicture}

{\small Una \omterm{def:g1:addition:def}{suma} en dibujos: los dos grupos juntos forman un solo
grupo de cinco.}
\end{omfigure}

\begin{example}[El orden no importa]\label{ex:g1:addition:order}
$2 + 3$ y $3 + 2$ juntan las mismas canicas: las dos dan $5$.
Para calcular $2 + 9$ es más rápido empezar por el número mayor:
di «nueve» y cuenta dos más --- «diez, once». Así,
$2 + 9 = 11$.
\end{example}

\section{Los amigos del diez}

\begin{example}[Complementos a 10]\label{ex:g1:addition:friends}
Dos números que juntos dan $10$ son «amigos del diez»:
\[
1 + 9, \quad 2 + 8, \quad 3 + 7, \quad 4 + 6, \quad 5 + 5 .
\]
Los diez dedos los enseñan todos: levanta $3$ dedos y los $7$ que quedan
doblados son su amigo. Saberse de memoria los amigos del diez hace que
muchas \omterm{def:g1:addition:def}{sumas} salgan solas.
\end{example}

\begin{omfigure}
\begin{tikzpicture}[scale=0.72]
  % tabla del diez: 2 filas de 5 casillas
  \draw[thick] (0,0) grid[step=1] (5,2);
  \draw[very thick] (0,0) rectangle (5,2);
  \foreach \i in {0,1,2} \fill[omDef] (\i+0.5,1.5) circle (0.3);
  \foreach \i in {3,4} \draw[omThm, thick] (\i+0.5,1.5) circle (0.3);
  \foreach \i in {0,...,4} \draw[omThm, thick] (\i+0.5,0.5) circle (0.3);
  \node[right, font=\small, align=left] at (5.6,1.0)
    {una \emph{tabla del diez}: $10$ casillas.\\
     $3$ llenas, $7$ vacías: $3 + 7 = 10$};
\end{tikzpicture}

{\small La tabla del diez enseña de un vistazo todas las parejas de amigos
del diez: lo que está lleno y lo que falta suman siempre $10$.}
\end{omfigure}

\begin{method}[Pasar por el diez]\label{met:g1:addition:crossten}
Para calcular $8 + 5$:
\begin{enumerate}
  \item quita del $5$ lo que le falta al $8$ para llegar a $10$: es
        $2$ (su amigo del diez);
  \item $8 + 2 = 10$, y del $5$ quedan $3$;
  \item $10 + 3 = 13$. Por tanto, $8 + 5 = 13$.
\end{enumerate}
Pasar por el $10$ convierte una \omterm{def:g1:addition:def}{suma} difícil en dos fáciles.
\end{method}

\begin{omfigure}
\begin{tikzpicture}[scale=0.62]
  % tabla del diez con 8 azules y 2 rojas (tomadas del 5)
  \draw[thick] (0,0) grid[step=1] (5,2);
  \draw[very thick] (0,0) rectangle (5,2);
  \foreach \i in {0,...,4} \fill[omDef] (\i+0.5,1.5) circle (0.3);
  \foreach \i in {0,1,2} \fill[omDef] (\i+0.5,0.5) circle (0.3);
  \foreach \i in {3,4} \fill[omThm] (\i+0.5,0.5) circle (0.3);
  % quedan 3 rojas fuera de la tabla
  \foreach \i in {0,1,2} \fill[omThm] (6.2+\i,1.0) circle (0.3);
  \node[below, font=\small] at (2.5,-0.25) {la tabla está llena: $10$};
  \node[below, font=\small] at (7.2,0.55) {sobran $3$};
\end{tikzpicture}

{\small $8 + 5$ en una tabla del diez: las $8$ fichas azules cogen $2$ de las
$5$ rojas para llenar la tabla --- una decena completa y $3$ más, así que
$13$.}
\end{omfigure}

\begin{omfigure}
\begin{tikzpicture}[scale=0.6]
  \draw[-Stealth] (-0.4,0) -- (15.4,0);
  \foreach \i in {0,...,15} \draw (\i,-0.1) -- (\i,0.1);
  \foreach \i in {0,5,8,10,13,15}
    \node[below, font=\footnotesize] at (\i,-0.12) {\i};
  \draw[-Stealth, omDef, thick] (8,0.3) .. controls (9,0.95) ..
    (10,0.3);
  \node[omDef, font=\small] at (9,1.05) {$+2$};
  \draw[-Stealth, omThm, thick] (10,0.3) .. controls (11.5,1.15) ..
    (13,0.3);
  \node[omThm, font=\small] at (11.5,1.2) {$+3$};
\end{tikzpicture}

{\small $8 + 5$ en la recta numérica: primero saltamos al $10$ y luego damos el resto.}
\end{omfigure}

\begin{example}[Los dobles]\label{ex:g1:addition:doubles}
Conviene aprenderse los dobles de memoria:
\[
1+1=2, \quad 2+2=4, \quad 3+3=6, \quad 4+4=8, \quad 5+5=10,
\]
\[
6+6=12, \quad 7+7=14, \quad 8+8=16, \quad 9+9=18, \quad 10+10=20 .
\]
Los casi dobles salen gratis: $6 + 7$ es uno más que $6 + 6$, o sea, $13$.
\end{example}

\section{Problemas pequeños}

\begin{example}\label{ex:g1:addition:problem}
«Lea tiene $7$ pegatinas. Le dan $6$ más. ¿Cuántas tiene ahora?»
Recibir más es sumar: $7 + 6 = 13$ (pasando por el diez:
$7 + 3 = 10$ y después $+3$). Lea tiene $13$ pegatinas. ¡Termina siempre
con una frase de respuesta!
\end{example}

\section{Ejercicios}

\begin{exercise}[$\star$]\label{exo:g1:addition:1}
Dibuja las canicas y calcula: $4 + 3$; \quad $5 + 2$; \quad
$6 + 4$.
\end{exercise}

\begin{exercise}[$\star$]\label{exo:g1:addition:2}
Calcula empezando por el número mayor: $3 + 8$; \quad
$2 + 12$; \quad $4 + 9$.
\end{exercise}

\begin{exercise}[$\star$]\label{exo:g1:addition:3}
Di el amigo del diez de: $7$; \ $4$; \ $9$; \ $5$; \ $1$.
\end{exercise}

\begin{exercise}[$\star$]\label{exo:g1:addition:4}
Calcula pasando por el diez (\cref{met:g1:addition:crossten}):
$9 + 4$; \quad $8 + 7$; \quad $6 + 5$.
\end{exercise}

\begin{exercise}[$\star$]\label{exo:g1:addition:5}
Recita los dobles y calcula con ellos: $7 + 8$; \quad
$5 + 6$; \quad $9 + 8$.
\end{exercise}

\begin{exercise}[$\star$]\label{exo:g1:addition:6}
Completa los huecos:
\[
6 + \;?\; = 10, \qquad
\;?\; + 5 = 12, \qquad
10 + \;?\; = 17 .
\]
\end{exercise}

\begin{exercise}[$\star$]\label{exo:g1:addition:7}
«En el autobús van $8$ niños; suben $5$ adultos. ¿Cuántas
personas hay en el autobús?» Escribe la \omterm{def:g1:addition:def}{suma} y la frase de
respuesta.
\end{exercise}

\begin{exercise}[$\star$]\label{exo:g1:addition:8}
Tomás tiene $6$ coches rojos y $7$ azules. Mía tiene $10$ coches. ¿Quién
tiene más coches?
\end{exercise}

\begin{exercise}[$\star\star$]\label{exo:g1:addition:9}
Tres dados marcan $4$, $6$ y $5$. ¿Cuál es el total? (\omterm{def:g1:addition:def}{Suma} primero dos
dados --- ¡elige los dos que son amigos del diez!)
\end{exercise}
